% Generated by generate_appendix_task_results.py. Do not edit the tables manually.

\section{Task Information}
\label{sec:appendix-task-information}

\begingroup
\footnotesize
\setlength{\tabcolsep}{3pt}
\setlength{\LTleft}{0pt}
\setlength{\LTright}{0pt}
\setlength{\LTcapwidth}{\linewidth}
\rowcolors{2}{RcbAppendixRowShade}{white}
\begin{longtable}{@{}>{\raggedright\arraybackslash}p{0.17\linewidth}>{\raggedright\arraybackslash}p{0.43\linewidth}>{\raggedright\arraybackslash}p{0.35\linewidth}@{}}
\caption{{\small Task information for the 40 ResearchClawBench tasks. Each row corresponds to one task. The Data column lists the \texttt{name}, modality, and \texttt{description} fields from each task's \texttt{task\_info.json}.}}\label{tab:appendix-task-metadata}\\
\toprule
\textbf{Task} & \textbf{Data} & \textbf{Scientific Objective} \\
\midrule
\endfirsthead
\toprule
\textbf{Task} & \textbf{Data} & \textbf{Scientific Objective} \\
\midrule
\endhead
\bottomrule
\endfoot
\texttt{Astronomy\_\allowbreak{}000} & \begin{minipage}[t]{\linewidth}\raggedright\vspace{0pt}IRAS\_\allowbreak{}09149-\allowbreak{}6206\_\allowbreak{}samples.\allowbreak{}dat (feature data): Contains the posterior distribution samples for the mass and dimensionless spin parameter of the supermassive black hole IRAS 09149-\allowbreak{}6206, providing the fundamental observational input for the constraint analysis in the supermassive black hole regime.\allowbreak{}\par\smallskip M33\_\allowbreak{}X-\allowbreak{}7\_\allowbreak{}samples.\allowbreak{}dat (feature data): Contains the posterior distribution samples for the mass and dimensionless spin parameter of the stellar-\allowbreak{}mass black hole in the X-\allowbreak{}ray binary M33 X-\allowbreak{}7, serving as the primary dataset for demonstrating and applying the Bayesian constraint framework.\allowbreak{}\end{minipage} & To constrain the properties of ultralight bosons (ULBs) by developing and applying a novel Bayesian statistical framework. This framework translates the physics of black hole superradiance into a probabilistic model that ingests the full posterior distributions (not just point estimates) of black hole mass and spin measurements. The goal is to derive statistically rigorous upper limits on ULB masses and self-interaction coupling strengths, thereby using astrophysical data to probe fundamental particle physics. \\
\texttt{Astronomy\_\allowbreak{}001} & \begin{minipage}[t]{\linewidth}\raggedright\vspace{0pt}DESI\_\allowbreak{}EDE\_\allowbreak{}Repro\_\allowbreak{}Data.\allowbreak{}txt (structure data): This dataset contains the best-\allowbreak{}fit parameters with 1sigma errors for LambdaCDM, EDE, and w0wa models from Tables II/\allowbreak{}III of the paper, along with manually extracted DESI BAO and Union3 SNe data points from Figure 6, used to reproduce the paper's key parameter constraints and distance comparison results.\allowbreak{}\end{minipage} & The scientific goal of this paper is to investigate whether an early dark energy (EDE) model can alleviate the acoustic tension between measurements from the cosmic microwave background (CMB) and baryon acoustic oscillations (BAO). Inputs: BAO data from DESI DR2, CMB data from Planck and ACT (including temperature and polarization power spectra and lensing), and Union3 supernova data (in some analyses). Outputs: Constraints on cosmological parameters (m, H0, sigma8, etc.) from model fitting, comparison of goodness-of-fit (2) for LambdaCDM, EDE, and w0wa models, and posterior distributions of EDE parameters (f\_EDE, log10a\_c). The results show that EDE can partially relieve the tension but leads to different parameter shifts compared to late-time dark energy models. \\
\texttt{Astronomy\_\allowbreak{}002} & \begin{minipage}[t]{\linewidth}\raggedright\vspace{0pt}H0DN\_\allowbreak{}MinimalDataset.\allowbreak{}txt (feature data): A minimal dataset to reproduce the Hubble constant measurement using the generalized least squares framework of the Distance Network, including geometric anchors, primary distance indicator measurements, secondary calibrations, and Hubble flow observations.\allowbreak{}\end{minipage} & The scientific goal of this paper is to achieve a \textasciitilde{}1\% precision measurement of the Hubble constant \textbackslash{}(H\_0\textbackslash{}) by constructing a "Local Distance Network" that combines multiple distance indicators through a covarianceweighted approach, providing a robust consensus result to address the Hubble tension. Inputs include geometric anchors (Milky Way parallaxes, LMC/SMC detached eclipsing binaries, NGC4258 masers), primary distance indicators (Cepheids, TRGB, Miras, JAGB), secondary indicators (SNe Ia, SBF, SNe II, FP, TF), and Hubbleflow measurements. Outputs are a consensus value of \textbackslash{}(H\_0\textbackslash{}) (baseline \textbackslash{}(H\_0 = 73.50 \textbackslash{}pm 0.81 \textbackslash{} \textbackslash{}mathrm\{km\textbackslash{},s\textasciicircum{}\{-1\}\textbackslash{},Mpc\textasciicircum{}\{-1\}\}\textbackslash{})), results from various analysis variants, and comparisons with earlyuniverse (CMB) constraints, along with publicly released software and data products. \\
\texttt{Astronomy\_\allowbreak{}003} & \begin{minipage}[t]{\linewidth}\raggedright\vspace{0pt}fig6\_\allowbreak{}data.\allowbreak{}csv (feature data): This dataset contains synthetic waveform differences representing the mismatch between the two highest numerical resolutions used in the SXS binary black hole simulations, after minimal time and phase alignment.\allowbreak{} The file has a single column with 1500 entries, each corresponding to one simulation in the catalog.\allowbreak{} The values are drawn from a lognormal distribution with a median of approximately 4x10-\allowbreak{}4, matching the typical resolution error reported in the SXS collaboration's third catalog paper.\allowbreak{} The distribution spans roughly 10-\allowbreak{}6 to 0.\allowbreak{}5, with a long tail toward larger differences.\allowbreak{} In the paper, such data are used to assess the overall numerical uncertainty of the waveform catalog and to demonstrate that the majority of simulations achieve high accuracy.\allowbreak{}\par\smallskip fig7\_\allowbreak{}data.\allowbreak{}csv (feature data): This file provides synthetic waveform differences decomposed by spherical harmonic mode l, covering l=2 through l=8.\allowbreak{} It consists of 1500 rows (simulations) and 7 columns, where each column corresponds to a specific l value and contains the minimalalignment waveform difference for that mode alone.\allowbreak{} The data are generated such that the median difference increases with l (from about 3x10-\allowbreak{}4 at l=2 to a few times 10-\allowbreak{}3 at l=8), and the scatter also grows slightly for higher l.\allowbreak{} In the original SXS study, such modal error distributions are critical for understanding how waveform accuracy varies across different multipoles and for guiding the truncation of mode contributions in gravitationalwave models.\allowbreak{}\par\smallskip fig8\_\allowbreak{}data.\allowbreak{}csv (feature data): This dataset contrasts waveform differences arising from two extrapolationorder comparisons: N=2 vs N=3 and N=2 vs N=4.\allowbreak{} It contains 1200 rows and two columns;\allowbreak{} the first column stores the differences between extrapolation orders 2 and 3, the second column stores differences between orders 2 and 4.\allowbreak{} The synthetic values are drawn from lognormal distributions with medians of 2x10-\allowbreak{}5 (for N2 vs N3) and 5x10-\allowbreak{}5 (for N2 vs N4), reflecting the trend that higherorder extrapolation pairs yield larger discrepancies.\allowbreak{} In the SXS catalog paper, such comparisons are used to evaluate the convergence of the extrapolation procedure that extracts waveforms from finiteradius simulation data to infinite null infinity, an essential step for producing reliable templates for gravitationalwave astronomy.\allowbreak{}\end{minipage} & Input: Initial parameters of binary black hole systems, including mass ratio, spin vectors, orbital eccentricity, etc. Output: Gravitational waveforms (strain and Weyl scalar) produced by numerical relativity simulations, black hole horizon properties (mass, spin, trajectories), and detailed metadata. Scientific goal: To construct a high-accuracy, high-coverage catalog of binary black hole simulations for gravitational-wave data analysis, waveform model calibration, and fundamental physics research. \\
\texttt{Chemistry\_\allowbreak{}000} & \begin{minipage}[t]{\linewidth}\raggedright\vspace{0pt}bace.\allowbreak{}csv (structure data): The BACE dataset contains small-\allowbreak{}molecule compounds represented by SMILES strings along with binary labels indicating whether each molecule inhibits human beta-\allowbreak{}secretase 1 (BACE-\allowbreak{}1).\allowbreak{} The dataset is used for molecular property prediction tasks in drug discovery, where molecular structures are converted into graph representations (atoms as nodes and bonds as edges) for classification modeling.\allowbreak{}\par\smallskip bbbp.\allowbreak{}csv (structure data): The BBBP (Blood-\allowbreak{}Brain Barrier Penetration) dataset contains small-\allowbreak{}molecule compounds represented by SMILES strings and binary labels indicating whether a compound can penetrate the blood-\allowbreak{}brain barrier.\allowbreak{} The dataset is used for molecular property classification tasks in pharmacology and drug design.\allowbreak{}\par\smallskip clintox.\allowbreak{}csv (structure data): The ClinTox dataset consists of small-\allowbreak{}molecule compounds represented by SMILES strings and labeled according to their clinical trial toxicity outcomes and FDA approval status.\allowbreak{} It is a multi-\allowbreak{}task binary classification dataset designed to evaluate a model's ability to predict both drug toxicity and regulatory approval likelihood.\allowbreak{} Molecules are converted into graph representations for molecular property prediction tasks.\allowbreak{}\par\smallskip hiv.\allowbreak{}csv (structure data): The HIV dataset contains small-\allowbreak{}molecule compounds represented by SMILES strings and labeled according to their ability to inhibit HIV replication.\allowbreak{} It is a binary classification dataset commonly used in molecular property prediction benchmarks.\allowbreak{} Each molecule is transformed into a graph structure for deep learning-\allowbreak{}based prediction of antiviral activity.\allowbreak{}\par\smallskip muv.\allowbreak{}csv (structure data): The MUV (Maximum Unbiased Validation) dataset is a large-\allowbreak{}scale molecular benchmark dataset consisting of small-\allowbreak{}molecule compounds represented by SMILES strings and labeled across multiple virtual screening tasks.\allowbreak{} It is designed to provide challenging and unbiased evaluation settings for molecular property prediction models.\allowbreak{} The dataset includes multiple binary classification tasks and exhibits high class imbalance, making it particularly difficult for predictive modeling.\allowbreak{}\end{minipage} & The task of this study is to design and evaluate a novel graph neural network architecture, termed Kolmogorov-Arnold Graph Neural Networks (KA-GNNs), for molecular property prediction by representing molecules as graphs with atom-level and bond-level features (including both covalent and non-covalent interactions) as input, and producing predictions of molecular properties such as toxicity, bioactivity, and physiological effects as output; the overarching scientific objective is to enhance predictive accuracy, computational efficiency, and interpretability by replacing conventional MLP-based transformations in graph neural networks with Fourier-based Kolmogorov-Arnold network modules that provide stronger expressive power and theoretical approximation guarantees. \\
\texttt{Chemistry\_\allowbreak{}001} & \begin{minipage}[t]{\linewidth}\raggedright\vspace{0pt}2l3r\_\allowbreak{}protein.\allowbreak{}pdb (protein structure): This file contains the experimental structure of the FKBP12 protein (PDB ID: 2L3R) determined by NMR.\allowbreak{} It includes only the CA atoms of all 107 residues in PDB format.\allowbreak{} This file serves as the ground truth protein structure for evaluating AlphaFold 3 predictions.\allowbreak{} The CA coordinates are directly comparable to the model's protein backbone output and are used in RMSD calculations and structural overlay visualizations.\allowbreak{}\par\smallskip 2l3r\_\allowbreak{}ligand.\allowbreak{}sdf (ligand structure): This file provides the experimental 3D conformation of the FK506 ligand, stored in the standard Structure-\allowbreak{}Data File (SDF) format.\allowbreak{} It contains the full atomic coordinates, bond connectivity, and chemical properties of the molecule.\allowbreak{} This is the primary reference for evaluating ligand pose prediction accuracy;\allowbreak{} the ligand RMSD is computed by aligning the predicted ligand coordinates to this reference using symmetry-\allowbreak{}aware Hungarian matching.\allowbreak{}\end{minipage} & Develop a unified deep learning framework that takes protein sequences, nucleic acid sequences, and small molecule structures as input, and outputs accurate 3D structures of biomolecular complexes using a diffusion-based architecture to predict interactions across diverse biological molecules. \\
\texttt{Chemistry\_\allowbreak{}002} & \begin{minipage}[t]{\linewidth}\raggedright\vspace{0pt}1brs\_\allowbreak{}AD.\allowbreak{}pdb (structure data): Processed structure of barnase-\allowbreak{}barstar complex (chains A and D) without water, used as input for HADDOCK3 analysis.\allowbreak{}\par\smallskip skempi\_\allowbreak{}v2.\allowbreak{}csv (feature data): SKEMPI 2.\allowbreak{}0 database containing experimental binding affinity changes upon mutation, used for validation.\allowbreak{}\end{minipage} & The input to HADDOCK3 consists of atomic coordinates of biomolecules (proteins, glycans, etc.) in PDB format, along with optional experimental restraints (e.g., ambiguous interaction restraints) and user-defined workflows. The output is an ensemble of modeled three-dimensional structures of biomolecular complexes, ranked and clustered according to various scoring functions. The scientific goal is to provide a versatile, modular platform for integrative modeling that leverages experimental data to predict accurate structures of biomolecular complexes, complementing machine learning approaches. \\
\texttt{Chemistry\_\allowbreak{}003} & \begin{minipage}[t]{\linewidth}\raggedright\vspace{0pt}random\_\allowbreak{}charges.\allowbreak{}xyz (structure data): This dataset contains 128-\allowbreak{}atom configurations with fixed point charges (+1e and -\allowbreak{}1e) randomly placed in a box.\allowbreak{} The interactions are modeled by Coulomb potential and a repulsive LennardJones term.\allowbreak{} It is used to benchmark whether the Latent Ewald Summation (LES) method can recover the exact atomic charges solely from energy and force data, as shown in Fig.\allowbreak{} 1 of the paper.\allowbreak{}\par\smallskip charged\_\allowbreak{}dimer.\allowbreak{}xyz (structure data): This dataset consists of configurations of two charged molecular dimers (each with total charges +1e and -\allowbreak{}1e) at various separation distances, with small internal distortions.\allowbreak{} It is employed to evaluate the ability of longrange models to capture binding energy curves when molecules are beyond the shortrange cutoff, as illustrated in Fig.\allowbreak{} 3 of the paper.\allowbreak{}\par\smallskip ag3\_\allowbreak{}chargestates.\allowbreak{}xyz (structure data): This dataset includes Ag3 trimers in two different charge states (+1 and -\allowbreak{}1) with varying bond lengths and random distortions.\allowbreak{} It is used to demonstrate that a shortrange model with global charge embedding (or separate training) is necessary to distinguish potential energy surfaces of different charge states, as shown in Fig.\allowbreak{} 5e and Table 1 of the paper.\allowbreak{}\end{minipage} & To develop a machine-learning interatomic potential that accurately and efficiently incorporates long-range electrostatic interactions without explicitly learning atomic charges or performing charge equilibration, thereby improving predictions for systems where electrostatics are critical (e.g., electrochemical interfaces, charged molecules, ionic liquids). \\
\texttt{Earth\_\allowbreak{}000} & \begin{minipage}[t]{\linewidth}\raggedright\vspace{0pt}glambie (sequence data): This dataset is the result of collecting, homogenizing, combining, and analyzing regional estimates derived from four primary observation methods (in situ glaciological measurements, digital elevation model differencing, satellite altimetry, and gravimetry) as well as integrated methods.\allowbreak{} It incorporates 233 regional estimates of glacier mass change contributed by 35 research teams and approximately 450 data contributors.\allowbreak{}\end{minipage} & Reconcile diverse observational methods to deliver a consistent and high-confidence assessment of global glacial mass change, and establish an observational benchmark for IPCC reports and climate model calibration. \\
\texttt{Earth\_\allowbreak{}001} & \begin{minipage}[t]{\linewidth}\raggedright\vspace{0pt}cloud\_\allowbreak{}seeding\_\allowbreak{}us\_\allowbreak{}2000\_\allowbreak{}2025 (sequence data): Official project-\allowbreak{}level cloud-\allowbreak{}seeding records released with the target paper, covering reported U.\allowbreak{}S.\allowbreak{} activities from 2000 to 2025.\allowbreak{} This is the only dataset used in this submission and it supports all reproduced tables, figures, and summary conclusions.\allowbreak{} The dataset contains 12 structured fields per record: filename, project name, year, season, state, operator affiliation, seeding agent, deployment apparatus, stated purpose, target area, control area, start date, and end date.\allowbreak{} The data enables comprehensive analysis of temporal trends, geographic distributions, operational characteristics, and methodological patterns in U.\allowbreak{}S.\allowbreak{} weather modification activities over a 25-\allowbreak{}year period.\allowbreak{}\end{minipage} & Input: NOAA weather-modification records released by the target paper, covering reported cloud-seeding projects in the United States from 2000 to 2025. Output: reproducible tables and figure-level evidence for spatial concentration, annual activity dynamics, purpose composition, and agent-apparatus deployment patterns. Scientific objective: test whether the paper's central empirical conclusions can be independently recovered from the published structured dataset using transparent, script-based analysis. \\
\texttt{Earth\_\allowbreak{}002} & \begin{minipage}[t]{\linewidth}\raggedright\vspace{0pt}gmw\_\allowbreak{}v4\_\allowbreak{}ref\_\allowbreak{}smpls\_\allowbreak{}qad\_\allowbreak{}v12.\allowbreak{}gpkg (vector data): Global mangrove extent polygons from Global Mangrove Watch (Bunting et al.\allowbreak{}, 2018), used to derive centroid points and calculate mangrove area.\allowbreak{} Sampled to 10\% for efficiency.\allowbreak{}\par\smallskip total\_\allowbreak{}ssp245\_\allowbreak{}medium\_\allowbreak{}confidence\_\allowbreak{}rates.\allowbreak{}nc (netcdf): Regional relative sea level rise rates for SSP2-\allowbreak{}4.\allowbreak{}5 (medium confidence) from IPCC AR6 (Garner et al.\allowbreak{}, 2021), used to extract median rates 2020-\allowbreak{}2100.\allowbreak{}\par\smallskip total\_\allowbreak{}ssp370\_\allowbreak{}medium\_\allowbreak{}confidence\_\allowbreak{}rates.\allowbreak{}nc (netcdf): Regional relative sea level rise rates for SSP3-\allowbreak{}7.\allowbreak{}0 (medium confidence) from IPCC AR6 (Garner et al.\allowbreak{}, 2021), used to extract median rates 2020-\allowbreak{}2100.\allowbreak{}\par\smallskip total\_\allowbreak{}ssp585\_\allowbreak{}medium\_\allowbreak{}confidence\_\allowbreak{}rates.\allowbreak{}nc (netcdf): Regional relative sea level rise rates for SSP5-\allowbreak{}8.\allowbreak{}5 (medium confidence) from IPCC AR6 (Garner et al.\allowbreak{}, 2021), used to extract median rates 2020-\allowbreak{}2100.\allowbreak{}\par\smallskip tracks\_\allowbreak{}mit\_\allowbreak{}mpi-\allowbreak{}esm1-\allowbreak{}2-\allowbreak{}hr\_\allowbreak{}historical\_\allowbreak{}reduced.\allowbreak{}nc (netcdf): Historical tropical cyclone tracks from the MIT model (Emanuel et al.\allowbreak{}, 2006) downscaled from CMIP6 MPI-\allowbreak{}ESM1-\allowbreak{}2-\allowbreak{}HR, covering 1850-\allowbreak{}2014.\allowbreak{} Used to calculate baseline cyclone frequencies after filtering and downsampling.\allowbreak{}\end{minipage} & The task of this paper is to develop a composite risk index combining tropical cyclone regime shifts and sea level rise, and apply it globally to evaluate where and to what extent mangroves and their ecosystem services are at risk by the end of the century, in order to inform climate-adaptive conservation and management strategies. \\
\texttt{Earth\_\allowbreak{}003} & \begin{minipage}[t]{\linewidth}\raggedright\vspace{0pt}20231012-\allowbreak{}06\_\allowbreak{}input\_\allowbreak{}netcdf.\allowbreak{}nc (sequence data): Pre-\allowbreak{}processed input data (20231012-\allowbreak{}06\_\allowbreak{}input\_\allowbreak{}netcdf.\allowbreak{}nc) with shape (2, 70, 721, 1440) representing two consecutive 6-\allowbreak{}hour atmospheric states at 0.\allowbreak{}25 deg resolution with 70 variables/\allowbreak{}channels\par\smallskip 006.\allowbreak{}nc (sequence data): FuXi output forecasts (006.\allowbreak{}nc) at 6-\allowbreak{}hour intervals up to 15 days.\allowbreak{}\end{minipage} & Input: Global atmospheric reanalysis data (ERA5) at 0.25 deg resolution, including 5 upper-air variables (geopotential, temperature, u-wind, v-wind, relative humidity) at 13 pressure levels and 5 surface variables (2m temperature, 10m u-wind, 10m v-wind, mean sea level pressure, total precipitation), from two consecutive 6-hour time steps. Output: 15-day global weather forecasts at 6-hour temporal resolution. Scientific Goal: Develop a cascade machine learning forecasting system using three specialized U-Transformer models to mitigate forecast error accumulation and extend skillful weather prediction to 15 days, achieving performance comparable to the ECMWF ensemble mean. \\
\texttt{Energy\_\allowbreak{}000} & \begin{minipage}[t]{\linewidth}\raggedright\vspace{0pt}NASA PCoE Dataset Repository (structure data): Experimental aging data of 18650 Li-\allowbreak{}ion batteries provided by the NASA Prognostics Center of Excellence (PCoE).\allowbreak{} It includes voltage, current, and temperature profiles recorded during constant current (CC) discharge cycles at room temperature, used here for experimental validation of the identification algorithm.\allowbreak{}\par\smallskip CS2\_\allowbreak{}36 (sequence data): Cycle life test data for a Commercial NCM (Nickel Cobalt Manganese) 18650 cell provided by the University of Maryland CALCE Battery Research Group.\allowbreak{} The dataset features standard 1C constant current discharge curves, used as the primary reference for parameter identification.\allowbreak{}\par\smallskip Oxford Battery Degradation Dataset (feature data): Long-\allowbreak{}term battery degradation data provided by the Oxford Battery Intelligence Lab.\allowbreak{} It contains dynamic urban driving profiles (highly transient current loads) obtained from 740mAh pouch cells, utilized to validate the model's generalization ability under dynamic conditions.\allowbreak{}\end{minipage} & (Definition of input, output, and scientific goal)Text to copy:Input: Experimental macroscopic data (voltage, temperature, and capacity curves under discharge conditions) and a multi-parameter search space defined by Latin Hypercube Sampling (LHS).Output: A set of identified high-fidelity internal parameters (such as particle radius, reaction rates, and thermal coefficients) for the electrochemical-aging-thermal (ECAT) coupled model.Scientific Goal: To develop a rapid and accurate parameter identification framework (MMGA) that uses an Artificial Neural Network (ANN) meta-model to replace computationally expensive physical simulations, thereby solving the trade-off between model complexity and calculation efficiency for Lithium-ion battery digital twins. \\
\texttt{Energy\_\allowbreak{}001} & \begin{minipage}[t]{\linewidth}\raggedright\vspace{0pt}buses.\allowbreak{}csv (structure data): Defines the buses (nodes) of the power system, including bus name, nominal voltage, and carrier type.\allowbreak{} This information forms the foundation of the grid topology.\allowbreak{}\par\smallskip links.\allowbreak{}csv (structure data): Describes transmission lines (or links), including source bus, target bus, nominal power capacity, line length, and carrier type.\allowbreak{} Used to simulate grid transmission capabilities and constraints.\allowbreak{}\par\smallskip demand.\allowbreak{}csv (sequence data): Provides hourly active power demand (MW) at each bus for 168 hours (one week).\allowbreak{} This is the electricity demand time series that the system must satisfy.\allowbreak{}\par\smallskip generators.\allowbreak{}csv (structure data): Lists all generator units with attributes such as bus location, carrier type (e.\allowbreak{}g.\allowbreak{}, wind, gas, nuclear), rated capacity, and marginal cost.\allowbreak{} Defines the generation resources of the system.\allowbreak{}\par\smallskip wind\_\allowbreak{}cf.\allowbreak{}csv (sequence data): Contains hourly wind capacity factors (0\textasciitilde{}1) for each bus, reflecting the temporal variability of wind resources.\allowbreak{} Used to calculate the maximum available output of wind turbines.\allowbreak{}\par\smallskip storage.\allowbreak{}csv (structure data): Describes the parameters of storage units, including bus location, type (e.\allowbreak{}g.\allowbreak{}, pumped hydro), power capacity, energy capacity, and charge/\allowbreak{}discharge efficiency.\allowbreak{} Used to simulate the charging and discharging behavior of storage.\allowbreak{}\end{minipage} & Input: Historical and future energy system data for Great Britain, including network topology, generator capacities, demand profiles, renewable time series, fuel prices, and National Grid's Future Energy Scenarios (FES) up to 2050. Output: Optimal power dispatch (generation, storage, curtailment) and system costs under different scenarios, with high spatial (29-node or zonal) and temporal (hourly) resolution. Scientific objective: To provide a fully open-source, high-resolution model of the GB power system that enables transparent, reproducible analysis of future energy pathways, such as renewable integration, network constraints, and flexibility options. \\
\texttt{Energy\_\allowbreak{}002} & \begin{minipage}[t]{\linewidth}\raggedright\vspace{0pt}hex\_\allowbreak{}final\_\allowbreak{}NA\_\allowbreak{}min.\allowbreak{}csv (feature data): Simulated dataset for African hydrogen production sites, including latitude, longitude, PV potential, wind potential, and distances to road, grid, ocean, and water infrastructure.\allowbreak{} Used as input for LCOH calculations.\allowbreak{}\par\smallskip ne\_\allowbreak{}10m\_\allowbreak{}admin\_\allowbreak{}0\_\allowbreak{}countries.\allowbreak{}shp (vector data): Main shapefile containing country boundary geometries for Africa and the world at 1:10m scale.\allowbreak{} Used as basemap for spatial visualizations.\allowbreak{}\par\smallskip ne\_\allowbreak{}10m\_\allowbreak{}admin\_\allowbreak{}0\_\allowbreak{}countries.\allowbreak{}shx (vector data): Shape index file required for reading the shapefile geometry data efficiently.\allowbreak{}\par\smallskip ne\_\allowbreak{}10m\_\allowbreak{}admin\_\allowbreak{}0\_\allowbreak{}countries.\allowbreak{}dbf (vector data): Attribute database file containing tabular information (country names, codes, etc.\allowbreak{}) associated with each geometry.\allowbreak{}\par\smallskip ne\_\allowbreak{}10m\_\allowbreak{}admin\_\allowbreak{}0\_\allowbreak{}countries.\allowbreak{}prj (vector data): Projection file that defines the coordinate system and map projection of the shapefile data.\allowbreak{}\par\smallskip ne\_\allowbreak{}10m\_\allowbreak{}admin\_\allowbreak{}0\_\allowbreak{}countries.\allowbreak{}cpg (vector data): Code page file specifying the character encoding used in the .\allowbreak{}dbf attribute file (typically UTF-\allowbreak{}8 or Latin-\allowbreak{}1).\allowbreak{}\end{minipage} & Build a transparent geospatial levelized-cost model to estimate the delivered cost of African green hydrogen to Europe (via ammonia shipping and reconversion) by 2030 under multiple financing and policy scenarios, identify least-cost/competitive locations, and quantify how de-risking and the interest-rate environment change cost competitiveness relative to producing green hydrogen in Europe. \\
\texttt{Energy\_\allowbreak{}003} & \begin{minipage}[t]{\linewidth}\raggedright\vspace{0pt}HEEW\_\allowbreak{}Mini-\allowbreak{}Dataset (sequence data): This compact and small dataset version of the core features of the target literature HEEW contains hourly electricity, heat, cooling load, photovoltaic power generation, greenhouse gas emissions, and seven meteorological data for the entire year of 2014.\allowbreak{} The data is organized in a hierarchical structure, covering 10 independent buildings (BN001-\allowbreak{}BN010), one aggregated community (CN01), and the entire area (Total), and is used to replicate the core experiments in the paper such as data cleaning, correlation analysis, and consistency verification of hierarchical aggregation, providing example data support for multi-\allowbreak{}energy system research and the development of machine learning algorithms.\allowbreak{}\end{minipage} & Input: Raw data sourced from sensor measurements (electricity, heat, cooling loads, PV generation of 147 buildings) of the Arizona State University Campus Metabolism Project and meteorological observations (temperature, humidity, wind speed, pressure, precipitation) from the U.S. National Weather Service. Output: A multi-source, hierarchical time-series dataset (HEEW) comprising 11,987,328 records with 13 hourly variables (electricity, heat, cooling loads, PV generation, GHG emissions, and 7 weather attributes) from 2014 to 2022, along with data cleaning algorithms. Scientific Goal: To provide a publicly available, comprehensive, and hierarchical benchmark dataset for energy system management, machine learning, and data-driven optimization (e.g., load forecasting, anomaly detection, clustering, imputation), addressing the gaps in existing multi-energy datasets regarding thermal loads, PV generation, emissions, and long-term coverage. \\
\texttt{Information\_\allowbreak{}000} & \begin{minipage}[t]{\linewidth}\raggedright\vspace{0pt}equation.\allowbreak{}png (sequence data): An image containing a mathematical equation used to evaluate the model's optical character recognition (OCR) and formula-\allowbreak{}to-\allowbreak{}LaTeX conversion capabilities.\allowbreak{}\par\smallskip doge.\allowbreak{}jpg (sequence data): A specific meme image ("Swole Doge vs.\allowbreak{} Cheems") used in Figure 5 of the paper.\allowbreak{} It contains embedded text ("Decoupling Visual Encoding" vs.\allowbreak{} "Single Visual Encoder") and visual metaphors to evaluate the model's high-\allowbreak{}level semantic understanding of humor.\allowbreak{}\end{minipage} & Build a unified autoregressive framework that decouples visual encoding to perform both multimodal understanding (e.g., visual question answering) and visual generation (e.g., text-to-image generation) within a single Transformer architecture. \\
\texttt{Information\_\allowbreak{}001} & \begin{minipage}[t]{\linewidth}\raggedright\vspace{0pt}demo\_\allowbreak{}imgs (sequence data): Two Demo pictures used in experiment 1\end{minipage} & The paper introduces a training-free framework designed to improve the fine-grained perception of MLLMs. The Scientific Objective is to mitigate the information loss caused by fixed-resolution vision encoders (like CLIP) when processing small objects. By using a task-guided cropping strategy, the model autonomously identifies regions of interest, 'zooms' into them, and integrates this local detail back into the global context to generate more accurate visual reasoning. \\
\texttt{Information\_\allowbreak{}002} & \begin{minipage}[t]{\linewidth}\raggedright\vspace{0pt}2111.\allowbreak{}01152 (feature data): The target scientific paper defining the AB-\allowbreak{}stacked MoTe2/\allowbreak{}WSe2 moire system and its Hamiltonian parameters.\allowbreak{}\end{minipage} & Input multi-step analytic calculation tasks of the Hartree-Fock method from 15 quantum many-body physics research papers; output correctly derived Hartree-Fock Hamiltonians, calculation step scores, and automated results of paper information extraction and step scoring; the scientific goal is to verify whether large language models (LLMs) can accurately perform research-level theoretical physics calculations via structured prompt templates and mitigate key bottlenecks in the research process. \\
\texttt{Information\_\allowbreak{}003} & \begin{minipage}[t]{\linewidth}\raggedright\vspace{0pt}NF-\allowbreak{}UNSW-\allowbreak{}NB15-\allowbreak{}v2\_\allowbreak{}3d.\allowbreak{}pt (feature data): NF-\allowbreak{}UNSW-\allowbreak{}NB15 is a NetFlowbased feature dataset where each row represents a single network flow described by 8 to 53 statistical features (e.\allowbreak{}g.\allowbreak{}, timestamp, duration, bytes, packet rates, interarrival times), extracted from packet headers and stored in CSV format for binary/\allowbreak{}multiclass intrusion detection.\allowbreak{}\end{minipage} & To address the inconsistent performance and poor detection capability of existing Network Intrusion Detection Systems (NIDS) across different attack types-especially for unknown and few-shot attacks-by proposing a disentangled dynamic intrusion detection framework (DIDS-MFL). The framework aims to disentangle entangled feature distributions in traffic data through statistical and representational disentanglement, incorporate dynamic graph diffusion for spatiotemporal aggregation, and enhance few-shot learning via multi-scale representation fusion, thereby improving detection accuracy, consistency, and generalization in real-world network environments. \\
\texttt{Life\_\allowbreak{}000} & \begin{minipage}[t]{\linewidth}\raggedright\vspace{0pt}Initial\_\allowbreak{}Training\_\allowbreak{}Data\_\allowbreak{}180 (feature data): Batch 1, The initial experimental dataset containing monomer compositions and adhesive strengths for the 180 bio-\allowbreak{}inspired hydrogels used to train the base models.\allowbreak{}\par\smallskip Initial\_\allowbreak{}Training\_\allowbreak{}Data\_\allowbreak{}180 (feature data): Batch 2, The initial experimental dataset containing monomer compositions and adhesive strengths for the 180 bio-\allowbreak{}inspired hydrogels used to train the base models.\allowbreak{}\par\smallskip Initial\_\allowbreak{}Training\_\allowbreak{}Data\_\allowbreak{}180 (feature data): Batch 3, The initial experimental dataset containing monomer compositions and adhesive strengths for the 180 bio-\allowbreak{}inspired hydrogels used to train the base models.\allowbreak{}\par\smallskip Initial\_\allowbreak{}Training\_\allowbreak{}Data (feature data): The cleaned and verified dataset containing the initial 184 hydrogel formulations.\allowbreak{} This file is the primary input for the 'rfr\_\allowbreak{}gp.\allowbreak{}py' script to train the initial machine learning models.\allowbreak{}\par\smallskip Final\_\allowbreak{}Optimization\_\allowbreak{}Dataset (feature data): The comprehensive dataset aggregating experimental results from all optimization rounds (1, 2, and 3).\allowbreak{} It serves as the input for evaluation notebooks to analyze the overall optimization trajectory and validate performance.\allowbreak{}\par\smallskip Final\_\allowbreak{}Optimization\_\allowbreak{}Dataset (feature data): The comprehensive dataset aggregating experimental results from another batch of final optimization dataset.\allowbreak{} It serves as the input for evaluation notebooks to analyze the overall optimization trajectory and validate performance.\allowbreak{}\end{minipage} & Input is protein sequence features converted to monomer compositions; Output is hydrogel adhesive strength. To de novo design synthetic hydrogels that achieve robust underwater adhesion (>1 MPa) by statistically replicating the sequence features of natural adhesive proteins. \\
\texttt{Life\_\allowbreak{}001} & \begin{minipage}[t]{\linewidth}\raggedright\vspace{0pt}cell-\allowbreak{}populations.\allowbreak{}csv (simulation output): Simulated cancer cell populations across repetitions;\allowbreak{} rows give cell ID, presented peptide, HLA allele, simulation name (e.\allowbreak{}g.\allowbreak{}, '100-\allowbreak{}cells.\allowbreak{}10x'), and mutation identifier.\allowbreak{}\par\smallskip final-\allowbreak{}response-\allowbreak{}likelihoods.\allowbreak{}csv (simulation output): Final immune response probability (p\_\allowbreak{}response), log value, presented-peptide count, population identifier, and vaccine type for each simulated cell;\allowbreak{} used for response distributions and coverage curves.\allowbreak{}\par\smallskip optimization\_\allowbreak{}runtime\_\allowbreak{}data.\allowbreak{}csv (performance metric): Optimization runtimes by population size and patient sample, used to generate Figure 6 (runtime vs population size).\allowbreak{}\par\smallskip selected-\allowbreak{}vaccine-\allowbreak{}elements.\allowbreak{}budget-\allowbreak{}10.\allowbreak{}minsum.\allowbreak{}adaptive.\allowbreak{}csv (optimization output): Lists the selected vaccine elements (peptides/\allowbreak{}mutations) under the MinSum objective with a budget of 10.\allowbreak{} Includes peptide, repetition, simulation name, weight, and run time.\allowbreak{} Used for recall analysis and vaccine composition.\allowbreak{}\par\smallskip sim-\allowbreak{}specific-\allowbreak{}response-\allowbreak{}likelihoods.\allowbreak{}csv (simulation output): Provides response probabilities for specific simulation repetitions.\allowbreak{} Similar to final-\allowbreak{}response-\allowbreak{}likelihoods but with more detailed simulation identifiers (including repetition).\allowbreak{}\par\smallskip vaccine-\allowbreak{}elements.\allowbreak{}scores.\allowbreak{}100-\allowbreak{}cells.\allowbreak{}10x.\allowbreak{}rep-\allowbreak{}0.\allowbreak{}csv (simulation output): One of 10 replicate files (rep-\allowbreak{}0 to rep-\allowbreak{}9) containing cell-\allowbreak{}level scores for each vaccine element.\allowbreak{} Each row gives the response probability for a specific cell and vaccine element, along with log probabilities.\allowbreak{} Used to aggregate cell response probabilities.\allowbreak{}\par\smallskip vaccine-\allowbreak{}elements.\allowbreak{}scores.\allowbreak{}100-\allowbreak{}cells.\allowbreak{}10x.\allowbreak{}rep-\allowbreak{}1.\allowbreak{}csv (simulation output): Replicate 1 of the vaccine element scores.\allowbreak{} Same structure as rep-\allowbreak{}0.\allowbreak{}\par\smallskip vaccine-\allowbreak{}elements.\allowbreak{}scores.\allowbreak{}100-\allowbreak{}cells.\allowbreak{}10x.\allowbreak{}rep-\allowbreak{}2.\allowbreak{}csv (simulation output): Replicate 2 of the vaccine element scores.\allowbreak{}\par\smallskip vaccine-\allowbreak{}elements.\allowbreak{}scores.\allowbreak{}100-\allowbreak{}cells.\allowbreak{}10x.\allowbreak{}rep-\allowbreak{}3.\allowbreak{}csv (simulation output): Replicate 3 of the vaccine element scores.\allowbreak{}\par\smallskip vaccine-\allowbreak{}elements.\allowbreak{}scores.\allowbreak{}100-\allowbreak{}cells.\allowbreak{}10x.\allowbreak{}rep-\allowbreak{}4.\allowbreak{}csv (simulation output): Replicate 4 of the vaccine element scores.\allowbreak{}\par\smallskip vaccine-\allowbreak{}elements.\allowbreak{}scores.\allowbreak{}100-\allowbreak{}cells.\allowbreak{}10x.\allowbreak{}rep-\allowbreak{}5.\allowbreak{}csv (simulation output): Replicate 5 of the vaccine element scores.\allowbreak{}\par\smallskip vaccine-\allowbreak{}elements.\allowbreak{}scores.\allowbreak{}100-\allowbreak{}cells.\allowbreak{}10x.\allowbreak{}rep-\allowbreak{}6.\allowbreak{}csv (simulation output): Replicate 6 of the vaccine element scores.\allowbreak{}\par\smallskip vaccine-\allowbreak{}elements.\allowbreak{}scores.\allowbreak{}100-\allowbreak{}cells.\allowbreak{}10x.\allowbreak{}rep-\allowbreak{}7.\allowbreak{}csv (simulation output): Replicate 7 of the vaccine element scores.\allowbreak{}\par\smallskip vaccine-\allowbreak{}elements.\allowbreak{}scores.\allowbreak{}100-\allowbreak{}cells.\allowbreak{}10x.\allowbreak{}rep-\allowbreak{}8.\allowbreak{}csv (simulation output): Replicate 8 of the vaccine element scores.\allowbreak{}\par\smallskip vaccine-\allowbreak{}elements.\allowbreak{}scores.\allowbreak{}100-\allowbreak{}cells.\allowbreak{}10x.\allowbreak{}rep-\allowbreak{}9.\allowbreak{}csv (simulation output): Replicate 9 of the vaccine element scores.\allowbreak{}\par\smallskip vaccine.\allowbreak{}budget-\allowbreak{}10.\allowbreak{}minsum.\allowbreak{}adaptive.\allowbreak{}csv (optimization output): Simplified vaccine composition file listing selected peptides and their counts/\allowbreak{}weights.\allowbreak{} Used for recall analysis.\allowbreak{}\end{minipage} & Input: Patient-specific sequencing data (tumor DNA/RNA, healthy DNA), HLA typing results, mutation VAF, gene expression (mean/variance), and prediction scores for peptide cleavage, MHC binding, and pMHC stability (from tools like pVACtools); vaccine manufacturing budget (maximum number of neoantigen elements). Output: Optimal personalized neoantigen vaccine composition (a set of neoantigen elements), quantitative vaccine efficacy metrics (per-cell immune response probability, coverage ratio of tumor cells, IoU of optimal vaccine compositions), and optimization runtime data. \\
\texttt{Life\_\allowbreak{}002} & \begin{minipage}[t]{\linewidth}\raggedright\vspace{0pt}7xg4.\allowbreak{}pdb (structure data): Query complex structure from PDB ID 7xg4 (Pseudomonas aeruginosa type IVA CRISPR-\allowbreak{}Cas system).\allowbreak{} Used in the paper as a known structural hit against an environmental Sulfitobacter sp.\allowbreak{} JL08 complex.\allowbreak{}\par\smallskip 6n40.\allowbreak{}pdb (structure data): Target complex structure from PDB ID 6n40.\allowbreak{} Used for pairwise structural alignment with 7xg4 to test FoldseekMultimer's alignment capability.\allowbreak{}\end{minipage} & The input is the three-dimensional structure of a protein complex (represented in PDB/mmCIF format or Foldseek database format), and the output is the structural alignment result between the complexes, including the correspondence between chains, superimposition vectors, and the TM score used to quantify similarity. Its scientific goal is to achieve efficient search and similarity detection in large-scale protein complex structure databases (such as those containing millions of structures) through ultra-fast and sensitive alignment algorithms. \\
\texttt{Life\_\allowbreak{}003} & \begin{minipage}[t]{\linewidth}\raggedright\vspace{0pt}dna\_\allowbreak{}r9.\allowbreak{}4.\allowbreak{}1\_\allowbreak{}400bps\_\allowbreak{}6mer\_\allowbreak{}uncalled4.\allowbreak{}csv (feature data): Contains 6-\allowbreak{}mer sequences and associated current statistics (mean, std, dwell time) for the DNA r9.\allowbreak{}4.\allowbreak{}1 pore model, used to generate substitution profiles and analyze base-\allowbreak{}position effects.\allowbreak{}\par\smallskip dna\_\allowbreak{}r10.\allowbreak{}4.\allowbreak{}1\_\allowbreak{}400bps\_\allowbreak{}9mer\_\allowbreak{}uncalled4.\allowbreak{}csv (feature data): Provides 9-\allowbreak{}mer pore model data for DNA r10.\allowbreak{}4.\allowbreak{}1 chemistry, including mean current, standard deviation, and dwell time, enabling comparison of substitution patterns between different pore versions.\allowbreak{}\par\smallskip rna\_\allowbreak{}r9.\allowbreak{}4.\allowbreak{}1\_\allowbreak{}70bps\_\allowbreak{}5mer\_\allowbreak{}uncalled4.\allowbreak{}csv (feature data): RNA001 pore model with 5-\allowbreak{}mer current parameters, used to study the relationship between nucleotide composition and ionic current in direct RNA sequencing.\allowbreak{}\par\smallskip rna004\_\allowbreak{}130bps\_\allowbreak{}9mer\_\allowbreak{}uncalled4.\allowbreak{}csv (feature data): RNA004 pore model (9-\allowbreak{}mer) containing current mean, standard deviation, and dwell time, facilitating comparison of RNA pore chemistries and their impact on signal characteristics.\allowbreak{}\par\smallskip performance\_\allowbreak{}summary.\allowbreak{}csv (feature data): Summary table of alignment time and file size for Uncalled4, f5c, Nanopolish, and Tombo across four sequencing chemistries, used to reproduce Table 1 performance benchmarks.\allowbreak{}\par\smallskip m6a\_\allowbreak{}predictions\_\allowbreak{}uncalled4.\allowbreak{}csv (feature data): m6Anet prediction probabilities for each candidate site based on Uncalled4 alignments, used to compute precision-\allowbreak{}recall curves and compare modification detection sensitivity.\allowbreak{}\par\smallskip m6a\_\allowbreak{}predictions\_\allowbreak{}nanopolish.\allowbreak{}csv (feature data): m6Anet prediction probabilities for the same sites using Nanopolish alignments, serving as a baseline for evaluating relative performance.\allowbreak{}\par\smallskip m6a\_\allowbreak{}labels.\allowbreak{}csv (feature data): Ground truth binary labels (0/\allowbreak{}1) for each site, derived from GLORI or m6A-\allowbreak{}Atlas, used as the reference for precision-\allowbreak{}recall and ROC analysis.\allowbreak{}\end{minipage} & To develop Uncalled4, a fast and accurate toolkit for nanopore signal alignment, enabling more sensitive and comprehensive detection of DNA and RNA modifications, while overcoming limitations of existing tools in speed, file format, and compatibility with new sequencing chemistries. \\
\texttt{Material\_\allowbreak{}000} & \begin{minipage}[t]{\linewidth}\raggedright\vspace{0pt}pretrain\_\allowbreak{}data.\allowbreak{}pt (structure data): This dataset contains a large number of unlabeled crystal structure graphs (5,000 samples) used for self-\allowbreak{}supervised pre-\allowbreak{}training.\allowbreak{} The model learns intrinsic representations of crystal structures without requiring magnetic labels, capturing general features that aid downstream tasks.\allowbreak{}\par\smallskip finetune\_\allowbreak{}data.\allowbreak{}pt (structure data): This labeled dataset (2,000 samples) consists of crystal graphs with binary labels indicating altermagnetic (positive) or non-\allowbreak{}altermagnetic (negative) materials.\allowbreak{} It simulates the scarcity of known altermagnets by having only 5\% positive samples (100 positives, 1900 negatives).\allowbreak{} It is used to fine-\allowbreak{}tune the pre-\allowbreak{}trained encoder into a classifier.\allowbreak{}\par\smallskip candidate\_\allowbreak{}data.\allowbreak{}pt (structure data): This unlabeled dataset (1,000 samples) represents candidate materials whose magnetic properties are unknown.\allowbreak{} The trained classifier predicts the probability of each being an altermagnet.\allowbreak{} Hidden true labels are stored internally for evaluation, allowing measurement of discovery accuracy.\allowbreak{} Approximately 50 true positives are embedded based on generation rules.\allowbreak{}\end{minipage} & The scientific objective of this work is to develop an AI-powered search engine that accelerates the discovery of new altermagnetic materials with targeted physical properties. The input consists of crystal structure data (represented as graphs) from databases such as the Materials Project, including a large unlabeled set for pre-training and a small labeled set of known altermagnets (148 positive samples) for fine-tuning. The output is a list of candidate materials predicted to be altermagnets with high probability, along with their electronic structure properties confirmed by first-principles calculations (e.g., 50 newly discovered altermagnets with classifications such as metal/insulator and d/g/i-wave anisotropy). \\
\texttt{Material\_\allowbreak{}001} & \begin{minipage}[t]{\linewidth}\raggedright\vspace{0pt}M-\allowbreak{}AI-\allowbreak{}Synth\_\allowbreak{}\_\allowbreak{}Materials\_\allowbreak{}AI\_\allowbreak{}Dataset\_\allowbreak{}.\allowbreak{}txt (structure data): This dataset is designed for rapid validation of three core AI application workflows in materials science-\allowbreak{}property prediction, structure generation, and experimental optimization-\allowbreak{}supporting code prototyping and fundamental algorithm testing.\allowbreak{}\end{minipage} & To accelerate the discovery, development, and optimization of advanced materials by integrating multimodal data through AI/ML models, enabling data-driven inverse design and reducing reliance on traditional trial-and-error approaches. \\
\texttt{Material\_\allowbreak{}002} & \begin{minipage}[t]{\linewidth}\raggedright\vspace{0pt}MACE-\allowbreak{}MP-\allowbreak{}0\_\allowbreak{}Reproduction Dataset.\allowbreak{}txt (structure data): This dataset contains all parameters and structural information needed to reproduce the performance of the MACE-\allowbreak{}MP-\allowbreak{}0 foundation model on three key tests: liquid water structure, adsorption energy scaling relations on transition metal surfaces, and CRBH20 reaction barriers.\allowbreak{}\end{minipage} & Input: The MPtrj dataset from the Materials Project (\textasciitilde{}1.5 million inorganic crystal structures and relaxation trajectories) and the MACE graph neural network architecture. Output: A general-purpose foundation model for atomistic potentials that can be directly applied to diverse chemical systems (liquids, solids, catalysis, reactions, etc.), with the ability to achieve ab initio accuracy after fine-tuning on minimal task-specific data. Scientific Goal: To develop and validate a universal foundation model for atomistic simulations that covers the periodic table, stably simulates diverse material systems, and achieves quantitative accuracy with minimal fine-tuning. \\
\texttt{Material\_\allowbreak{}003} & \begin{minipage}[t]{\linewidth}\raggedright\vspace{0pt}tg\_\allowbreak{}calibration (feature data): Contains molecular SMILES, experimental Tg values, and MD simulated Tg values used to train and evaluate the Gaussian process calibration model.\allowbreak{}\par\smallskip tg\_\allowbreak{}vitrimer\_\allowbreak{}MD (feature data): Contains molecular structures of vitrimer systems and their MD simulated Tg values used as input for the Gaussian process calibration to generate calibrated Tg predictions.\allowbreak{}\end{minipage} & Develop an AI-guided inverse-design framework for recyclable vitrimeric polymers by combining molecular dynamics simulations, Gaussian-process calibration, and a graph variational autoencoder, with the goal of generating new vitrimer chemistries that achieve desired glass transition temperatures (Tg) and validating selected candidates experimentally. \\
\texttt{Math\_\allowbreak{}000} & \begin{minipage}[t]{\linewidth}\raggedright\vspace{0pt}simulated\_\allowbreak{}sequence.\allowbreak{}json (structured data): This dataset contains a simulated multi-\allowbreak{}object video sequence generated with controlled parameters (40 frames, 20 objects, 85\% detection rate, 20\% occlusion overlap threshold) to evaluate tracking performance under dense occlusion scenarios.\allowbreak{} It includes ground truth trajectories and detection boxes with confidence scores and occlusion labels, enabling reproducible comparison of SparseTrack and ByteTrack as presented in the paper.\allowbreak{}\end{minipage} & input: Consecutive image frames and object detections per frame (including bounding box coordinates and confidence scores). outpt: Complete trajectories for each target, i.e., identity labels (IDs) and corresponding bounding box sequences across video frames. scientific targets: To effectively handle occlusions and improve multi-object tracking performance in crowded scenes by decomposing dense target sets into sparse subsets via pseudo-depth estimation and performing hierarchical association. \\
\texttt{Math\_\allowbreak{}001} & \begin{minipage}[t]{\linewidth}\raggedright\vspace{0pt}complex\_\allowbreak{}optimization\_\allowbreak{}data.\allowbreak{}npy (structure data): A synthetic, ill-\allowbreak{}conditioned dataset generated for high-\allowbreak{}dimensional Lasso regression.\allowbreak{} It contains the design matrix A (1000x2000), response vector b, and ground truth sparse coefficients x\_\allowbreak{}true, stored in .\allowbreak{}npy format to verify algorithm convergence.\allowbreak{}\end{minipage} & Input: A smooth convex objective function f(x) (potentially with a non-smooth regularization term) and an initial starting point x\_0. Output: The optimal solution x* that minimizes the global objective function with an accelerated convergence rate. Scientific Goal: To establish a unified Variable and Operator Splitting (VOS) framework that derives Nesterov's accelerated method and ADMM from a continuous-time dynamical system perspective, proving linear convergence using strong Lyapunov functions. \\
\texttt{Math\_\allowbreak{}002} & \begin{minipage}[t]{\linewidth}\raggedright\vspace{0pt}maps\_\allowbreak{}60\_\allowbreak{}10\_\allowbreak{}10\_\allowbreak{}0.\allowbreak{}175 (structure data): The task set consists of static 2D grid maps )and multi-\allowbreak{}agent task configurations defined by preset parameters such as agent count, start positions, and target positions.\allowbreak{}\par\smallskip empty (structure data): Dataset of empty 2D grid maps (likely 25x25) with no static obstacles, containing multi-\allowbreak{}agent task configurations to evaluate navigation in high-\allowbreak{}density open spaces without structural blockages.\allowbreak{}\par\smallskip maze (structure data): Dataset of maze-\allowbreak{}structured 2D grid maps (25x25) characterized by complex corridors and dead-\allowbreak{}ends, containing task configurations designed to test pathfinding algorithms in highly constrained environments.\allowbreak{}\par\smallskip random\_\allowbreak{}large (structure data): Dataset of large-\allowbreak{}scale (50x50) 2D grid maps with randomly generated obstacles (17.\allowbreak{}5\% density), containing task configurations serving as a benchmark for algorithm scalability in unstructured environments.\allowbreak{}\par\smallskip random\_\allowbreak{}medium (structure data): Dataset of medium-\allowbreak{}sized (25x25) 2D grid maps with randomly distributed obstacles (17.\allowbreak{}5\% density), containing task configurations that balance map complexity and computational load for standard evaluation.\allowbreak{}\par\smallskip random\_\allowbreak{}small (structure data): Dataset of small-\allowbreak{}scale (10x10) 2D grid maps with random obstacles (17.\allowbreak{}5\% density), containing task configurations primarily used for rapid testing and analyzing algorithm behavior in tight spaces.\allowbreak{}\par\smallskip room (structure data): Dataset of room-\allowbreak{}structured 2D grid maps (25x25) simulating indoor environments with connected chambers and narrow doorways, containing task configurations for analyzing bottleneck traversal.\allowbreak{}\par\smallskip warehouse (structure data): Dataset of warehouse-\allowbreak{}style 2D grid maps (25x25) with organized shelf layouts, containing task configurations specifically designed to simulate automated logistics and retrieval scenarios.\allowbreak{}\end{minipage} & Input: A MAPF (Multi-Agent Path Finding) instance I, which consists of a discrete 2D grid map with static obstacles and a set of agents, each with a distinct start position and a designated goal position. Output: A solution path set P, consisting of collision-free paths for all agents that navigate them from their start positions to their goal positions without vertex or swapping collisions. Scientific Goal (Task): To solve the Multi-Agent Path Finding problem by developing a hybrid algorithm that integrates Multi-Agent Reinforcement Learning into the Large Neighborhood Search framework. The specific objective is to balance solution quality (reducing collisions via MARL in early stages) and computational efficiency (using Prioritized Planning in later stages) to achieve higher success rates in complex environments compared to existing methods. \\
\texttt{Math\_\allowbreak{}003} & \begin{minipage}[t]{\linewidth}\raggedright\vspace{0pt}imo\_\allowbreak{}ag\_\allowbreak{}30.\allowbreak{}txt (structure data): A curated benchmark of 30 geometry problems from the International Mathematical Olympiad (since 2000), used for final evaluation.\allowbreak{}\end{minipage} & Input: Formal statements of olympiad-level geometry problems (e.g., IMO diagrams and premises). Output: Machine-verifiable, human-readable proofs for Euclidean geometry theorems. Scientific Goal: To develop an AI system that autonomously solves complex geometry problems without human demonstrations, advancing neuro-symbolic reasoning in mathematics. \\
\texttt{Neuroscience\_\allowbreak{}000} & \begin{minipage}[t]{\linewidth}\raggedright\vspace{0pt}Together\_\allowbreak{}1\_\allowbreak{}features\_\allowbreak{}extracted.\allowbreak{}csv (feature data): Frame-\allowbreak{}level engineered features derived from tracked animal pose signals in the official SimBA sample project;\allowbreak{} used as model input matrix X for both behavior labels.\allowbreak{}\par\smallskip Together\_\allowbreak{}1\_\allowbreak{}targets\_\allowbreak{}inserted.\allowbreak{}csv (sequence data): Frame-\allowbreak{}aligned target annotations for Attack and Sniffing from the same sample project;\allowbreak{} used as supervised targets y during train/\allowbreak{}test evaluation.\allowbreak{}\par\smallskip Together\_\allowbreak{}1\_\allowbreak{}machine\_\allowbreak{}results\_\allowbreak{}reference.\allowbreak{}csv (feature data): Reference output table provided with the sample project, retained as auxiliary material for contextual comparison with reproduced classifier outputs.\allowbreak{}\end{minipage} & Input: pose-derived frame-level feature tables and aligned behavior labels (Attack, Sniffing) from the official SimBA sample project. Output: trained supervised classifiers, quantitative evaluation reports, precision-recall diagnostics, confusion matrices, and feature-importance tables. Scientific objective: verify, on open data and executable code, whether the SimBA-style workflow can reproducibly transform tracked behavior features into transparent and auditable behavior classification evidence. \\
\texttt{Neuroscience\_\allowbreak{}001} & \begin{minipage}[t]{\linewidth}\raggedright\vspace{0pt}flow (sequence data): the complete ensemble of 50 pre-\allowbreak{}trained deep mechanistic network (DMN) models, along with all necessary configuration files, synapse count matrices, and celltype annotations.\allowbreak{} These models are constrained by the fly connectome and optimized for optic flow estimation, allowing users to directly simulate neural responses, reproduce key analyses from the paper, and generate experimentally testable hypotheses without retraining.\allowbreak{}\end{minipage} & Input The connectome of the motion pathway in the Drosophila optic lobe (including the number of synaptic connections, polarity, and spatial distribution of 64 cell types). Unknown single-neuron kinetic parameters (time constants, resting potentials) and unit synaptic strengths, which need to be determined through optimization. Output A deep mechanistic network (DMN) whose structure strictly follows the connectome, with parameters learned through task optimization, capable of simulating the voltage activities of 45,669 neurons in response to visual stimuli. Detailed predictions of the neural activity of each neuron, as well as quantitative analysis of motion detection mechanisms. Scientific Goal To demonstrate that the activity of each neuron in a neural circuit can be accurately predicted solely based on connectome measurements (structure) and task knowledge (functional goals), thereby establishing a bridge from structure to function. Overall Task Construct a connectome-constrained and task-optimized deep mechanistic network that can perform optical flow estimation tasks, and reveal the computational role of each neuron in the Drosophila visual system in motion detection through model predictions. \\
\texttt{Neuroscience\_\allowbreak{}002} & \begin{minipage}[t]{\linewidth}\raggedright\vspace{0pt}test\_\allowbreak{}simulated.\allowbreak{}csv (structure data): Contains approximately 3600 samples (30\% of total).\allowbreak{} Identical structure to the training set: 20 features, label, and degradation type.\allowbreak{} Used for evaluating model performance on unseen data.\allowbreak{}\par\smallskip train\_\allowbreak{}simulated.\allowbreak{}csv (structure data): Contains approximately 8400 samples (70\% of total).\allowbreak{} Each sample has 20 feature columns (019) representing morphology, intensity, and embedding modalities, a binary label (1 for same neuron, 0 otherwise), and a degradation type (Misalignment, Missing Sections, Mixed, or Average).\allowbreak{} The data is stratified by degradation to ensure balanced representation.\allowbreak{}\end{minipage} & Input: An over-segmented electron microscopy (EM) image volume of a fly brain and a pair of adjacent neuron segments (a query segment and a candidate segment) located near a potential truncation point Output: A binary prediction (0 or 1) indicating whether the two given segments belong to the same neuron and should be merged. Scientific Goal: To automate the proofreading process in large-scale connectomics by accurately predicting connectivity between over-segmented neuron fragments, thereby reducing the massive manual workload required to reconstruct complete neurons from petascale EM data. \\
\texttt{Neuroscience\_\allowbreak{}003} & \begin{minipage}[t]{\linewidth}\raggedright\vspace{0pt}adata\_\allowbreak{}RPE.\allowbreak{}h5ad (feature data): A preprocessed single-\allowbreak{}cell dataset (protein iterative indirect immunofluorescence imaging) showcasing cellular state transitions in a retina-\allowbreak{}related context that is compatible with neuroscience-\allowbreak{}adjacent analysis.\allowbreak{}\end{minipage} & The input is single-cell readouts (such as scRNA-seq or protein imaging), and the output is a selected subset of dynamically expressed molecular features that best preserves continuous cellular trajectories. This setup supports analyses of neural lineage progression, glial activation, and neurodegeneration-related state transitions while reducing confounding variation. \\
\texttt{Physics\_\allowbreak{}000} & \begin{minipage}[t]{\linewidth}\raggedright\vspace{0pt}Multi-\allowbreak{}component Icosahedral Reproduction Data.\allowbreak{}txt (sequence data): This dataset contains complete parameters and result data for reproducing all simulation experiments from the paper "General theory for packing icosahedral shells into multi-\allowbreak{}component aggregates", enabling full reproduction of theoretical calculations, experimental verification, and dynamic growth simulations in any computational environment.\allowbreak{}\end{minipage} & To establish a universal theoretical framework for the rational design of multi-component nanoclusters and nanoparticles with specific symmetry (chiral or achiral) and compositional sequences, and to predict their stability and growth behavior, ultimately enabling targeted material fabrication for applications in catalysis, optics, and related fields. \\
\texttt{Physics\_\allowbreak{}001} & \begin{minipage}[t]{\linewidth}\raggedright\vspace{0pt}MATBG Superfluid Stiffness Core Dataset.\allowbreak{}txt (feature data): This dataset fully contains all simulated data required to reproduce the three core experiments of the target study, covering carrier density dependence, temperature dependence, and current dependence.\allowbreak{} It can be directly used to independently verify key conclusions such as quantum geometry-\allowbreak{}dominated enhancement of superfluid stiffness, power-\allowbreak{}law behavior of anisotropic gaps, and quadratic current relationships.\allowbreak{}\end{minipage} & Input A magic-angle twisted bilayer graphene (MATBG) device with gate-tunable carrier density, subjected to DC bias current and microwave probe signals at cryogenic temperatures (\textasciitilde{}20 mK). Output The device's DC resistance, microwave resonance frequency, and their dependence on temperature, gate voltage, and current. The core extracted physical quantity is the superfluid stiffness and its temperature and current dependence. Scientific Goal To directly measure the superfluid stiffness of MATBG, test whether it significantly exceeds predictions of conventional Fermi liquid theory, investigate its power-law temperature dependence to reveal the nature of unconventional pairing (anisotropic gap), and verify the crucial role of quantum geometric effects in flat-band superconductivity. \\
\texttt{Physics\_\allowbreak{}002} & \begin{minipage}[t]{\linewidth}\raggedright\vspace{0pt}results (sequence data): The measurement result subset output from the Random Quantum Circuit Sampling (RCS) experiment is saved per circuit instance as results/\allowbreak{}N40\_\allowbreak{}verification/\allowbreak{}N40\_\allowbreak{}d*\allowbreak{}\_\allowbreak{}XEB/\allowbreak{}*\allowbreak{}\_\allowbreak{}counts.\allowbreak{}json.\allowbreak{} Each JSON file records several measured bitstrings (represented as tuple-\allowbreak{}strings or bitstrings) and their occurrence counts, which are used to compute the counts-\allowbreak{}weighted XEB fidelity estimate.\allowbreak{}\par\smallskip amplitudes (sequence data): The corresponding subset of ideal distribution information is saved per circuit instance as amplitudes/\allowbreak{}N40\_\allowbreak{}verification/\allowbreak{}N40\_\allowbreak{}d*\allowbreak{}\_\allowbreak{}XEB/\allowbreak{}*\allowbreak{}\_\allowbreak{}amplitudes.\allowbreak{}json.\allowbreak{} Each JSON file provides the ideal amplitudes or ideal probabilities for a set of bitstrings (unified as ideal\_\allowbreak{}probs in your code), which are used to match with the measured counts and compute XEB (the number of matched keys is typically 20 in your reproduction).\allowbreak{}\end{minipage} & Evaluation of the computational power of random quantum circuit sampling (RCS) on arbitrary geometries as presented in the paper .Input: Experimental sampling results (bitstring counts/samples) and their corresponding ideal distribution information (which can be the full ideal probability/amplitude or that of a verifiable subset) for different qubit counts N, circuit depths d, and instance indices r. The data should enable the implementation of the fidelity estimation workflow used in the paper (e.g., XEB, MB regression probability, gatecount/error propagation models, etc.). Output: A fidelity estimate with uncertainty for each (N, d, r) configuration. Under settings such as scanning d for fixed N or scanning N for fixed d, comparative curves should be plotted to validate the paper's core conclusion regarding the "gap between the experimental fidelity and classical approximability under arbitrarygeometry/highconnectivity random circuits." \\
\texttt{Physics\_\allowbreak{}003} & \begin{minipage}[t]{\linewidth}\raggedright\vspace{0pt}raw\_\allowbreak{}trARPES\_\allowbreak{}data.\allowbreak{}h5 (structure data): Raw, unprocessed 4D data arrays (energy, momentum kx/\allowbreak{}ky, time delay) from the tr-\allowbreak{}ARPES experiment.\allowbreak{}\par\smallskip processed\_\allowbreak{}band\_\allowbreak{}data.\allowbreak{}json (feature data): Processed data containing the extracted positions and intensities of the main Dirac cone and replica bands.\allowbreak{}\par\smallskip polarization\_\allowbreak{}dependence\_\allowbreak{}data.\allowbreak{}csv (sequence data): Tabular data containing the measured intensity of the replica band for each pump polarization angle (thetap).\allowbreak{}\end{minipage} & Input: Monolayer epitaxial graphene samples and mid-infrared pump excitation parameters (wavelength: 5 m, intensity, polarization angle). Output: Direct, energy- and momentum-resolved observation of Floquet-Bloch states (replica bands of the Dirac cone) via time-resolved and angle-resolved photoemission spectroscopy (tr-ARPES). Scientific Goal: To experimentally confirm the existence of Floquet-Bloch states in a paradigmatic 2D material and elucidate the underlying scattering mechanism involving photon-dressed Volkov final states. \\
\end{longtable}
\rowcolors{1}{white}{white}
\endgroup
\section{Per-Task Results}
\label{sec:appendix-per-task-results}

\begingroup
\scriptsize
\captionsetup{justification=raggedright,singlelinecheck=false}
\setlength{\tabcolsep}{2pt}
\begin{longtable}{@{}>{\raggedright\arraybackslash}p{0.170\linewidth}*{8}{>{\centering\arraybackslash}p{0.094\linewidth}}@{}}
\caption{{\small Per-task total rubric scores for each task and system. Panel (a) reports autonomous-agent scores: C.Code: Claude Code; Codex: Codex CLI; ARIS: ARIS Codex; Open: OpenClaw; Nano: Nanobot; Evo0: EvoScientist v0.0.4; Evo1: EvoScientist v0.1.1; RClaw: ResearchClaw. Panel (b) reports ResearchHarness LLM scores: C4.6/C4.7: Claude-Opus-4.6/4.7; DS: DeepSeek-V4-Pro; GLM: GLM-5.1; G5.4/G5.5: GPT-5.4/5.5; Gem/GemF: Gemini-3.1-Pro/Gemini-3.5-Flash; G4.1/G4.3: Grok-4.1/4.3; K2.5/K2.6: Kimi-K2.5/2.6; M2P/M2.5: MiMo-V2-Pro/V2.5; Q3.5/Q3.6/Q3.7: Qwen3.5-397B-A17B/Qwen3.6-Plus/Qwen3.7-Max. A dash indicates that the model directly crashed or failed during execution.}}\label{tab:appendix-task-scores}\label{tab:appendix-agent-scores}\label{tab:appendix-llm-scores}\\
\toprule
\multicolumn{9}{@{}l}{\textbf{(a) Autonomous agents}}\\
\textbf{Task} & \textbf{C.Code} & \textbf{Codex} & \textbf{ARIS} & \textbf{Open} & \textbf{Nano} & \textbf{Evo0} & \textbf{Evo1} & \textbf{RClaw} \\
\midrule
\endfirsthead
\toprule
\multicolumn{9}{@{}l}{\textbf{Table~\ref{tab:appendix-task-scores} continued. Panel (a) Autonomous agents.}}\\
\textbf{Task} & \textbf{C.Code} & \textbf{Codex} & \textbf{ARIS} & \textbf{Open} & \textbf{Nano} & \textbf{Evo0} & \textbf{Evo1} & \textbf{RClaw} \\
\midrule
\endhead
\bottomrule
\endfoot
\texttt{Astronomy\_\allowbreak{}000} & 23.5 & 29.4 & 16.5 & 18.6 & 14.5 & 11.5 & 24 & 20.8 \\
\texttt{Astronomy\_\allowbreak{}001} & 27.4 & 20 & 10.4 & 36 & 20.4 & 21.2 & 12.8 & 14.4 \\
\texttt{Astronomy\_\allowbreak{}002} & 23.5 & 12.7 & 11 & 11.6 & 10 & 23.9 & 23.1 & 11.1 \\
\texttt{Astronomy\_\allowbreak{}003} & 46.5 & 44 & 47.4 & 47.3 & 44.2 & 47.3 & 45.6 & 45.9 \\
\texttt{Chemistry\_\allowbreak{}000} & 14.75 & 11 & 15.4 & 12.85 & 16.85 & 7.5 & 9.55 & 18.25 \\
\texttt{Chemistry\_\allowbreak{}001} & 3.5 & 6.2 & 2.5 & 0.5 & 0.5 & 7.3 & 11.3 & 2.15 \\
\texttt{Chemistry\_\allowbreak{}002} & 1 & 1 & 0 & 1 & 0 & 1.4 & 3 & 0 \\
\texttt{Chemistry\_\allowbreak{}003} & 18 & 12.3 & 11.7 & 9.5 & 7.5 & 1.5 & 15.9 & 13.5 \\
\texttt{Earth\_\allowbreak{}000} & 24.6 & 25.3 & 14.5 & 20.1 & 21.5 & 12.8 & 23 & 14.7 \\
\texttt{Earth\_\allowbreak{}001} & 33.56 & 40.87 & 34.48 & 31.38 & 25.3 & 39.97 & 39.36 & 32.97 \\
\texttt{Earth\_\allowbreak{}002} & 31.7 & 22.6 & 11.4 & 16.8 & 9.9 & 15.1 & 13 & 20.6 \\
\texttt{Earth\_\allowbreak{}003} & 1 & 0 & 0 & 0.9 & 0.6 & 0 & 0 & 0 \\
\texttt{Energy\_\allowbreak{}000} & 11.3 & 13.2 & 10.7 & 16 & 9.5 & 12.4 & 9.9 & 3.5 \\
\texttt{Energy\_\allowbreak{}001} & 21.7 & 24.5 & 3.6 & 15.5 & 7.6 & 20.9 & 16.6 & 15.4 \\
\texttt{Energy\_\allowbreak{}002} & 38.6 & 33 & 32.55 & 31 & 27.85 & 36.4 & 33.6 & 37.65 \\
\texttt{Energy\_\allowbreak{}003} & 15.3 & 21.5 & 7.5 & 25.5 & 9 & 26.3 & 19.5 & 19.5 \\
\texttt{Information\_\allowbreak{}000} & 48 & 10 & 1.2 & 35.1 & 22.7 & 7.2 & 24.5 & 10 \\
\texttt{Information\_\allowbreak{}001} & 7 & 5.6 & 3.6 & 0 & 1 & 3.6 & 2.4 & 0 \\
\texttt{Information\_\allowbreak{}002} & 27.1 & 36.8 & 11.4 & 15.8 & 11.5 & 14.1 & 35 & 39.8 \\
\texttt{Information\_\allowbreak{}003} & 17.75 & 15.6 & 7.8 & 5.2 & 9.2 & 4.75 & 13.6 & 9.8 \\
\texttt{Life\_\allowbreak{}000} & 8.2 & 7.55 & 5.95 & 5.4 & 7.95 & 7.05 & 9.5 & 10.05 \\
\texttt{Life\_\allowbreak{}001} & 19.55 & 12.7 & 14 & 13.35 & 9.65 & 13.6 & 5.95 & 5.7 \\
\texttt{Life\_\allowbreak{}002} & 6.9 & 1.5 & 7.3 & 9.7 & 4.9 & 9.1 & 11.2 & 6.6 \\
\texttt{Life\_\allowbreak{}003} & 27.6 & 36 & 40.5 & 34.5 & 29.3 & 36 & 40.1 & 33.6 \\
\texttt{Material\_\allowbreak{}000} & 14.5 & 17.5 & 11.5 & 12 & 23 & 17.5 & 5.6 & 12.1 \\
\texttt{Material\_\allowbreak{}001} & 23.6 & 11.55 & 8.5 & 13.45 & 1.75 & 6.3 & 3.5 & 12.95 \\
\texttt{Material\_\allowbreak{}002} & 35.1 & 0.99 & 11.75 & 6.12 & 14.9 & 10.05 & 37.11 & 35.75 \\
\texttt{Material\_\allowbreak{}003} & 28.8 & 21.8 & 18 & 20.1 & 12.35 & 19.95 & 14.4 & 16.25 \\
\texttt{Math\_\allowbreak{}000} & 26.2 & 26.65 & 7.8 & 24.7 & 17.7 & 8 & 22.2 & 7.2 \\
\texttt{Math\_\allowbreak{}001} & 44.1 & 39 & 31.3 & 23.9 & 30.6 & 35.2 & 24.1 & 37.7 \\
\texttt{Math\_\allowbreak{}002} & 10.2 & 7.4 & 6.5 & 8.6 & 0 & 14.1 & 9.6 & 6.3 \\
\texttt{Math\_\allowbreak{}003} & 29.6 & 10 & 0 & 0 & 0 & 0 & 6 & 0 \\
\texttt{Neuroscience\_\allowbreak{}000} & 8.6 & 14 & 5 & 13.2 & 7.4 & 13.6 & 14.2 & 7.8 \\
\texttt{Neuroscience\_\allowbreak{}001} & 4.95 & 2.55 & 2 & 3.75 & 1.2 & 0.75 & 18.5 & 0.5 \\
\texttt{Neuroscience\_\allowbreak{}002} & 0.75 & 2 & 1 & 1 & 0.4 & 1 & 0.75 & 1 \\
\texttt{Neuroscience\_\allowbreak{}003} & 7.7 & 15.8 & 19.65 & 15.05 & 4.25 & 5.95 & 12.3 & 7.35 \\
\texttt{Physics\_\allowbreak{}000} & 27.4 & 32.1 & 18.2 & 7.5 & 7.8 & 19.2 & 24.4 & 22.7 \\
\texttt{Physics\_\allowbreak{}001} & 30.25 & 25.7 & 25.9 & 32 & 27.45 & 24.5 & 12.4 & 29.25 \\
\texttt{Physics\_\allowbreak{}002} & 25.05 & 21.4 & 20 & 27.45 & 10.35 & 24.85 & 47.4 & 24.25 \\
\texttt{Physics\_\allowbreak{}003} & 46.5 & 45 & 34.6 & 43.3 & 31.8 & 36.8 & 55.6 & 44.3 \\

\end{longtable}
\endgroup

\begingroup
\tiny
\setlength{\tabcolsep}{1pt}
\renewcommand{\arraystretch}{1.05}
\begin{longtable}{@{}>{\raggedright\arraybackslash}p{0.145\linewidth}*{17}{>{\centering\arraybackslash}p{0.047\linewidth}}@{}}
\toprule
\multicolumn{18}{@{}l}{\textbf{(b) ResearchHarness LLMs}}\\
\textbf{Task} & \textbf{C4.6} & \textbf{C4.7} & \textbf{DS} & \textbf{GLM} & \textbf{G5.4} & \textbf{G5.5} & \textbf{Gem} & \textbf{GemF} & \textbf{G4.1} & \textbf{G4.3} & \textbf{K2.5} & \textbf{K2.6} & \textbf{M2P} & \textbf{M2.5} & \textbf{Q3.5} & \textbf{Q3.6} & \textbf{Q3.7} \\
\midrule
\endfirsthead
\toprule
\multicolumn{18}{@{}l}{\textbf{Table~\ref{tab:appendix-task-scores} continued. Panel (b) ResearchHarness LLMs.}}\\
\textbf{Task} & \textbf{C4.6} & \textbf{C4.7} & \textbf{DS} & \textbf{GLM} & \textbf{G5.4} & \textbf{G5.5} & \textbf{Gem} & \textbf{GemF} & \textbf{G4.1} & \textbf{G4.3} & \textbf{K2.5} & \textbf{K2.6} & \textbf{M2P} & \textbf{M2.5} & \textbf{Q3.5} & \textbf{Q3.6} & \textbf{Q3.7} \\
\midrule
\endhead
\bottomrule
\endfoot
\texttt{Astronomy\_\allowbreak{}000} & 33.1 & 25.6 & 25.6 & 9 & 24.9 & 26 & 7.5 & 24 & 20 & -- & 14.4 & 3.1 & 11.2 & 19 & 11.1 & 21 & 15.9 \\
\texttt{Astronomy\_\allowbreak{}001} & -- & 26.6 & 6 & 37 & 4 & 2 & 4 & 13.2 & 29.4 & 7.2 & 21.2 & 36 & 10 & 8 & 8 & 35.2 & 11 \\
\texttt{Astronomy\_\allowbreak{}002} & 24.5 & 32 & 22.1 & 29.9 & 14.6 & 24.5 & 20 & -- & 20.4 & 20.6 & 11 & 28.1 & 19.1 & 27.5 & 3.1 & 20 & 28 \\
\texttt{Astronomy\_\allowbreak{}003} & 47.6 & 47.4 & 47.1 & 43.1 & 44.8 & 43.4 & 45.8 & 47 & 44.7 & 46.4 & 46.8 & 44 & 46.7 & 43.1 & 46.4 & 46.9 & 40.2 \\
\texttt{Chemistry\_\allowbreak{}000} & 9.35 & -- & 15.1 & 18.95 & 17.05 & 19.1 & 19.6 & 9 & 3.8 & 13.3 & 18.9 & -- & 14 & -- & 20.4 & 17.9 & 13.25 \\
\texttt{Chemistry\_\allowbreak{}001} & 9 & 7.4 & 5 & 2.75 & 0.5 & 2 & 4.25 & 6.5 & 1.5 & 0.5 & 3.55 & 5.75 & 0.75 & 5.25 & 0.5 & 3.95 & 0.5 \\
\texttt{Chemistry\_\allowbreak{}002} & 3.4 & 1 & 0 & 1.4 & 1 & 5 & 0 & 1 & 0.4 & 0 & 1.2 & 0.4 & 1 & 0 & 3.6 & 0 & 1 \\
\texttt{Chemistry\_\allowbreak{}003} & 7.2 & -- & 10 & 22.5 & 9 & 15.5 & -- & -- & 6 & 0.3 & 6.9 & -- & 9.5 & 8 & 6.9 & 9 & 12.9 \\
\texttt{Earth\_\allowbreak{}000} & 21.5 & 15.5 & 16.7 & 14.8 & 12.1 & 21.8 & 14.8 & 25.5 & 11.9 & 11.9 & 15.8 & 21.6 & 14.7 & 15.2 & 15.6 & 14.2 & 12.5 \\
\texttt{Earth\_\allowbreak{}001} & 40.94 & 38.94 & 37.38 & 33.04 & 36.71 & 39.24 & 28.05 & 33.19 & 30.78 & 35.56 & 21.06 & 41.65 & 35.57 & 22.54 & 35.58 & 26.63 & 36.65 \\
\texttt{Earth\_\allowbreak{}002} & 28.8 & 34 & 25.5 & 30.4 & 25 & 11.8 & 11 & 27 & 13.6 & 13 & 25.2 & 24.8 & 28.1 & 25.9 & 18.6 & 24.4 & 9.6 \\
\texttt{Earth\_\allowbreak{}003} & -- & 1.6 & 1 & 4 & 1.5 & 0.6 & 1.5 & 10.5 & 3.5 & 0 & 3.6 & 3.4 & 1 & 0 & 3.5 & 1.2 & 0.6 \\
\texttt{Energy\_\allowbreak{}000} & 12.3 & -- & 17.5 & 16 & 11.8 & 16 & 3.5 & 4.4 & 2.3 & 14.5 & 5 & 9.5 & 12.5 & 14.7 & 16.9 & 22 & -- \\
\texttt{Energy\_\allowbreak{}001} & 25.1 & 17.1 & 25 & 18.2 & 8.5 & 20.1 & 3 & 21 & 8 & -- & 3.4 & 10.4 & 9.1 & 23.4 & 10.6 & 24.7 & 27 \\
\texttt{Energy\_\allowbreak{}002} & 42.45 & 41 & 38.65 & 39.85 & 28.75 & 37 & 21.25 & 31.7 & 28.2 & 33.35 & 25.3 & 31.4 & 21.9 & 23.7 & 31.55 & 30.55 & 32.35 \\
\texttt{Energy\_\allowbreak{}003} & 24 & 11.5 & 6 & 7.5 & 19.5 & 17.5 & 20.3 & 16.3 & 6 & 8.7 & 19.5 & 21 & 7.5 & 16.3 & 15.5 & 7.5 & 16.5 \\
\texttt{Information\_\allowbreak{}000} & 23.2 & 0.8 & 0 & 19.5 & 29 & 22.7 & 10.5 & 49.4 & 25.5 & 2 & 16.4 & 0 & -- & 16.7 & 9.5 & 47.1 & -- \\
\texttt{Information\_\allowbreak{}001} & 3.6 & 7.6 & 1 & 0.6 & 0 & 6 & 1 & 9 & 5 & 3.6 & 7.6 & 18 & -- & 1 & 6.4 & 5.6 & -- \\
\texttt{Information\_\allowbreak{}002} & 20.1 & 32.3 & 12.9 & 36.2 & 22.9 & 33.1 & 8.8 & 27.2 & 10.8 & 1.8 & 10.3 & 27.8 & 13.2 & 26.9 & 7.2 & 12 & 38.4 \\
\texttt{Information\_\allowbreak{}003} & 4.35 & 14.9 & 4.6 & 8.65 & 11 & 9.55 & -- & 10.8 & 6.2 & 2.65 & 3.25 & 5.25 & 5.5 & 4.6 & 0.75 & 5.95 & 5.25 \\
\texttt{Life\_\allowbreak{}000} & 9.35 & 5.45 & 10.15 & 9.95 & 8.85 & 2.25 & 3.55 & 7.7 & 0.5 & 4.35 & 7.45 & 6.8 & 4.9 & 4.75 & 4.8 & 7 & -- \\
\texttt{Life\_\allowbreak{}001} & 3.6 & 12.45 & 9.5 & 8.75 & 7.45 & 7.2 & 8.5 & 12.95 & 13.7 & 11.05 & 7.75 & 11.55 & 7.05 & 11.95 & 10.75 & 5.7 & 0 \\
\texttt{Life\_\allowbreak{}002} & 5.7 & 6.5 & 1 & 6.3 & 8.8 & 8.5 & 7.5 & 6.3 & 2 & 5.5 & 5 & 2.8 & 6.6 & 3.5 & 6.5 & 7.3 & 6.1 \\
\texttt{Life\_\allowbreak{}003} & 32.1 & 26.9 & 32.5 & 23.8 & 31.8 & 29.4 & 18 & 26.5 & 36.5 & 29 & 25.4 & 34.1 & 32.6 & 30.9 & 23.8 & 28.4 & 24.8 \\
\texttt{Material\_\allowbreak{}000} & 25.1 & 21.6 & 24 & 21.7 & 5 & 5.6 & 20 & 22.6 & 16.4 & 20.9 & 4 & 15.5 & 13.2 & 23.5 & 14 & 20.8 & -- \\
\texttt{Material\_\allowbreak{}001} & 1.75 & 18 & 21 & 4.15 & 3.5 & 8 & 5.25 & 0 & 11.55 & 6.85 & 6.3 & 10.1 & 18 & 7 & 17.1 & 15.95 & 15.9 \\
\texttt{Material\_\allowbreak{}002} & 39.08 & 35.81 & 38.4 & 28.4 & 3.3 & 26.09 & -- & 40.71 & 0.66 & 16.98 & 31.15 & 39.08 & 37.31 & 28.5 & 33.5 & 27.05 & 30.49 \\
\texttt{Material\_\allowbreak{}003} & 14 & 21.15 & 15.15 & 20.3 & 18.8 & 18.95 & 8.7 & 15.6 & 9.2 & 4.65 & 10.95 & 24.65 & 17 & 17.7 & 8.9 & 14.45 & 25.25 \\
\texttt{Math\_\allowbreak{}000} & 23 & 22.95 & 23.35 & 23.95 & 11 & 8.25 & 21.1 & -- & 7.7 & 6.2 & 18.3 & 23.55 & 7.2 & 18.5 & 23.1 & 23.5 & 29.05 \\
\texttt{Math\_\allowbreak{}001} & 34.3 & 24.4 & 22.9 & 27 & 42 & 23 & 13.5 & 22 & 26.4 & 13.8 & 15.2 & 21.6 & 33.6 & 30.2 & 30.4 & 33.1 & 31.2 \\
\texttt{Math\_\allowbreak{}002} & 20.5 & 9.4 & 7.5 & 10.9 & 10.2 & 11 & -- & -- & 2.6 & 2 & 14.7 & 14 & 3.5 & 14.1 & 3.5 & 9.6 & 19.1 \\
\texttt{Math\_\allowbreak{}003} & 11.75 & -- & -- & 10 & 10 & 10 & 10.35 & -- & -- & 12 & -- & 18.75 & 10 & 10 & 0 & 10 & 10 \\
\texttt{Neuroscience\_\allowbreak{}000} & 10.4 & 10.6 & 10 & 10.4 & 13 & 11 & 2 & 8 & 8.4 & 6 & 7.4 & 9 & 9.2 & 3.6 & 11.4 & 8.6 & 11 \\
\texttt{Neuroscience\_\allowbreak{}001} & 2.7 & -- & 7.05 & 2.55 & 2.2 & 3.75 & -- & 0 & 4 & 1.2 & 0.75 & 3.5 & 2.7 & 1.25 & 1.25 & 2 & 3.75 \\
\texttt{Neuroscience\_\allowbreak{}002} & 0.75 & 3.75 & 0 & 1 & 0 & 1 & 0 & 0 & 0 & 0 & 0 & 0.5 & 0 & 1 & 0 & 0.75 & 0 \\
\texttt{Neuroscience\_\allowbreak{}003} & 15.55 & 14.7 & 13.65 & 9.45 & 4.35 & 9.6 & 2.7 & 0.75 & 6.75 & 3.2 & 3.7 & 20.4 & 13.05 & 12.75 & 4.05 & 7.2 & 13.15 \\
\texttt{Physics\_\allowbreak{}000} & 26.8 & 14.2 & 14.1 & 23.8 & 19.2 & 18.3 & 22 & 14.5 & 15.4 & 35.7 & 13.2 & 25.6 & 17.5 & 27.4 & 20.9 & 24.8 & -- \\
\texttt{Physics\_\allowbreak{}001} & -- & 39.95 & 38.6 & 31.9 & 28.5 & 29.75 & 22.55 & 31.65 & 34.7 & 17.95 & 32.05 & 17 & 22.55 & 46.25 & 25.8 & 27.55 & 29.45 \\
\texttt{Physics\_\allowbreak{}002} & 32.4 & 38.6 & 28.2 & 32.9 & 18.4 & 26.4 & 27.55 & -- & 11.25 & 19 & 20.1 & 21.1 & 22 & 29 & 19.55 & 28.5 & 40.05 \\
\texttt{Physics\_\allowbreak{}003} & 45.8 & 44.1 & 33.6 & 27.1 & 42.1 & 49 & 46.8 & 34.3 & 40.9 & 40 & 40.5 & 32.1 & 43.5 & 36 & 38.3 & 41.8 & 45.4 \\
\end{longtable}
\endgroup
